%%%%%%%%%%%%%%%%%%%%%%%%%%%%%%%%%%%%%%%%%%%%%%%%%%%%%%%%%%%%%%%%%%%%%%
% Lớp 10 | Véctơ (chưa xét tọa độ) | Tích vô hướng (chưa xét tọa độ)
% Dạng 0H5?4-3: Đẳng thức về tích vô hướng hoặc độ dài
% Số câu: 5  (sinh tự động bởi tach_cau.py; từ nay thầy sửa trực tiếp file này)
%%%%%%%%%%%%%%%%%%%%%%%%%%%%%%%%%%%%%%%%%%%%%%%%%%%%%%%%%%%%%%%%%%%%%%

%%%%%%%%%%%%%%%% TRẮC NGHIỆM 4 PHƯƠNG ÁN %%%%%%%%%%%%%%%%

% nguồn: MaHoaNH_TracNghiem_Lop10
\begin{ex}%[0H5V4-3]%[ma:e9a0b15c]
	Cho $\left|\overrightarrow{a}\right|=2$, $\left|\overrightarrow{b}\right|=3$, $\left|\overrightarrow{a}+2\overrightarrow{b}\right|=5$.  Tìm $\left|3\overrightarrow{a}-\overrightarrow{b}\right|$.
	\choice 
		 {$135$}
		 {$11$}
		{\True $\dfrac{3 \sqrt{30}}{2}$} 
		{$45$} 
	 \loigiai{
		Ta  có $\left|\overrightarrow{a}+2\overrightarrow{b}\right|^2=\left(\overrightarrow{a}+2\overrightarrow{b}\right)^2=\overrightarrow{a}^2+4\overrightarrow{a}\cdot\overrightarrow{b}+4\overrightarrow{b}^2=\left|\overrightarrow{a}\right|^2+4\overrightarrow{a}\cdot\overrightarrow{b}+4\left|\overrightarrow{b}\right|^2$.\\
		Suy ra \begin{align*}
		 \overrightarrow{a}\cdot\overrightarrow{b}=\dfrac{1}{4}\left(\left|\overrightarrow{a}+2\overrightarrow{b}\right|^2-\left|\overrightarrow{a}\right|^2-4\left|\overrightarrow{b}\right|^2\right)=\dfrac{1}{4}\left(5^2-2^2-4\cdot 3^2\right)=-\dfrac{15}{4}.
	\end{align*}
	Khi đó, \begin{align*}
	\left|3\overrightarrow{a}-\overrightarrow{b}\right|^2&=\left(3\overrightarrow{a}-\overrightarrow{b}\right)^2=9\overrightarrow{a}^2-6\overrightarrow{a}\cdot\overrightarrow{b}+\overrightarrow{b}^2\\ 
	 &=9\left|\overrightarrow{a}\right|^2-6\overrightarrow{a}\cdot\overrightarrow{b}+\left|\overrightarrow{b}\right|^2=9\cdot 2^2+\dfrac{45}{2}+3^2=\dfrac{135}{2}.
	\end{align*}
Vậy  $\left|3\overrightarrow{a}-\overrightarrow{b}\right|=\dfrac{3\sqrt{10}}{2}$.}
\end{ex}

% nguồn: MaHoaNH_TracNghiem_Lop10
\begin{ex}%[0H5V4-3]%[ma:3ddacfe7]
	Cho $\left|\overrightarrow{a}+\overrightarrow{b}\right|=4$, $\left|\overrightarrow{a}\right|=2$,  $\left|\overrightarrow{b}\right|=3$. Tính $\left|\overrightarrow{a}-\overrightarrow{b}\right|$.
	\choice 
		{$3$} 
		{\True  $\sqrt{10}$}
		 {$\sqrt{12}$}
		 {$2$}
	\loigiai{ 
		Ta có \begin{align*}
		\overrightarrow{a}\cdot\overrightarrow{b}&=\dfrac{1}{2}\left[\left(\overrightarrow{a}+\overrightarrow{b}\right)^2-\overrightarrow{a}^2-\overrightarrow{b}^2\right]=\dfrac{1}{2}\left(\left|\overrightarrow{a}+\overrightarrow{b}\right|^2-\left|\overrightarrow{a}\right|^2-\left|\overrightarrow{b}\right|^2\right)\\ 
		 &=\dfrac{1}{2}\left(4^2-2^2-3^2\right)=\dfrac{3}{2}.
	\end{align*}
	Khi đó, \begin{align*}
	\left|\overrightarrow{a}-\overrightarrow{b}\right|^2&=\left(\overrightarrow{a}-\overrightarrow{b}\right)^2=\overrightarrow{a}^2-2\overrightarrow{a}\cdot\overrightarrow{b}+\overrightarrow{b}^2\\ 
	&=\left|\overrightarrow{a}\right|^2-2\overrightarrow{a}\cdot\overrightarrow{b}+\left|\overrightarrow{b}\right|^2=2^2-3+3^2=10. 
	\end{align*}
	Vậy  $\left|\overrightarrow{a}-\overrightarrow{b}\right|=\sqrt{10}$.
} 
\end{ex}


%%%%%%%%%%%%%%%% TRẢ LỜI NGẮN %%%%%%%%%%%%%%%%

% nguồn: MaHoaNH_TraLoiNgan_Lop10
\begin{bt}%[0H5H4-3]%[ma:e57a7956]
	Cho hai vectơ $\overrightarrow{a}$ và $\overrightarrow{b}$. Biết $\left|\overrightarrow{a}\right|=2,\left|\overrightarrow{b}\right|=\sqrt{3}$ và $\left(\overrightarrow{a}, \overrightarrow{b}\right)=120^{\circ}$. Tính $\left|\overrightarrow{a}+\overrightarrow{b}\right|$ (làm tròn kết quả đến  hàng phần trăm).
	\shortans{1,88} 
	 \loigiai{
		Ta có \begin{align*}
		 \left|\overrightarrow{a}+\overrightarrow{b}\right|&=\sqrt{\left(\overrightarrow{a}+\overrightarrow{b}\right)^2}=\sqrt{\overrightarrow{a}^2+\overrightarrow{b}^2+2 \overrightarrow{a}\cdot\overrightarrow{b}}\\
		&=\sqrt{\left|\overrightarrow{a}\right|^2+\left|\overrightarrow{b}\right|^2+2\left|\overrightarrow{a}\right|  \cdot\left|\overrightarrow{b}\right| \cdot \cos \left(\overrightarrow{a}, \overrightarrow{b}\right)}=\sqrt{7-2 \sqrt{3}}\\
		&\approx 1{,}88. 
\end{align*}}
\end{bt}

% nguồn: MaHoaNH_TraLoiNgan_Lop10
\begin{bt}%[0H5H4-3]%[ma:04b2dfec]
	Cho tam giác $ABC$ có $BC=a$, $CA=b$, $AB=c$.  Biết $M$ là trung điểm của $BC$. Tính $\overrightarrow{AM}^2=\dfrac{m\left(b^2+c^2\right)+n\cdot a^2}{4}$. Khi đó $m-n$ bằng bao nhiêu?
	\shortans{3} 
	\loigiai{ 
	 \immini{
		Vì $M$ là  trung điểm của $BC$ nên
		$\overrightarrow{AM}=\dfrac{1}{2}\left(\overrightarrow{AB}+\overrightarrow{AC}\right)$.\\ 
		Khi  đó $\overrightarrow{AM}^2=\dfrac{1}{4}\left(\overrightarrow{AB}+\overrightarrow{AC}\right)^2=\dfrac{1}{4}\left(\overrightarrow{AB}^2+2 \overrightarrow{AB} \cdot \overrightarrow{AC}+\overrightarrow{AC}^2\right)$.\\
		Mà  $\overrightarrow{AB} \cdot \overrightarrow{AC}=\dfrac{b^2+c^2-a^2} 2$.\\
		Suy ra  $\overrightarrow{AM}^2=\dfrac{1}{4}\left(c^2+2 \cdot \dfrac{b^2+c^2-a^2} 2+b^2\right)=\dfrac{2\left(b^2+c^2\right)-a^2} 4$.\\
		Vậy  $m-n=2-\left(-1\right)=3$.\\
		(Đây cũng là  công thức để tính độ dài đường trung tuyến tam giác). }{
		\begin{tikzpicture}[line join = round,line cap = round,scale=0.3]
				\def \a{3}; 
				\coordinate  (B) at (0,0);
				\coordinate (C)  at ($(B)+(3*\a,0)$);
				\coordinate[shift=(60:(\a/2)  cm)] (A) at (B);
				\coordinate (M) at  ($(C)!0.5!(B)$);
				 \draw (C)--(A)--(B)--(C)--cycle (A)--(M);
				\draw[->,arrows={-latex}] (A)--(B); 
				\draw[->,arrows={-latex}]  (A)--(M);
				\draw[->,arrows={-latex}]  (A)--(C);
				 \foreach \x/\y in {A/90,B/-90,C/-90,M/-90}
				\draw[fill=black] (\x) circle  (1pt) + (\y:0.75cm) node{$\x$};
				 \foreach \x/\y in {B/M,M/C}{
					\path (\x)--(\y)  node[midway,sloped]{\tikz{\draw (-90:2pt)--(90:2pt);}};
					\draw (\x)--(\y) node[midway,sloped]  {};
				} 
			\end{tikzpicture}
		 }
	} 
\end{bt}


%%%%%%%%%%%%%%%% TỰ LUẬN %%%%%%%%%%%%%%%%

% nguồn: MaHoaNH_TuLuan_Lop10
\begin{bt}%[0H5C4-3]%[ma:b307ccd1]
	Cho tam giác $ABC$ có  $M$ là trung điểm $BC$ và $H$ là trực tâm. Biết $\overrightarrow{MH}\cdot \overrightarrow{MA}=k BC^2$. Tìm $k$.
	 \loigiai{
		Hình vẽ 
		\begin{center}
			\begin{tikzpicture}[scale=1,font=\footnotesize,line join=round,line cap=round,>=stealth]
					\path (0,0) coordinate  (B)
						(5,0) coordinate (C) 
						(1.5,3) coordinate (A) 
						($(B)!(C)!(A)$)  coordinate (D)
						($(A)!(B)!(C)$) coordinate (E) 
						(intersection of  C--D and B--E) coordinate (H)
						($(B)!0.5!(C)$)  coordinate (M);
					\draw  (A)--(B)--(C)--cycle (C)--(D)
						 (B)--(E)
						(A)--(M)--(H); 
					\foreach \p/\g in {A/90,B/-150,C/-30,D/180,E/30,H/90,M/-90} \fill[black]  (\p) circle (1pt)+(\g:0.3) node{$\p$};
					\foreach \x/\y/\z in {B/D/C,B/E/C} \draw pic[draw,angle radius=0.2cm]{right angle=\x--\y--\z}; 
				\end{tikzpicture}
		\end{center}
		 Vì $M$ là trung điểm $BC$ nên $\overrightarrow{AM}=\dfrac{1}{2}\left(\overrightarrow{AB}+\overrightarrow{AC}\right)$ và $\overrightarrow{H M}=\dfrac{1}{2}\left(\overrightarrow{HB}+\overrightarrow{HC}\right)$.\\
		Ta  có
		\begin{align*}
			\overrightarrow{MH}\cdot \overrightarrow{MA}&=\overrightarrow{AM}\cdot  \overrightarrow{H M}\\
			 &=\dfrac{1}{4}\left(\overrightarrow{AB}+\overrightarrow{AC}\right) \cdot\left(\overrightarrow{HB}+\overrightarrow{HC}\right)\\
			&=\dfrac{1}{4}\left(\overrightarrow{AB}\cdot \overrightarrow{HB}+\overrightarrow{AB}\cdot  \overrightarrow{HC}+\overrightarrow{AC}\cdot \overrightarrow{HB}+\overrightarrow{AC}\cdot \overrightarrow{HC}\right) \\
			 & =\dfrac{1}{4}\left(\overrightarrow{AB}\cdot \overrightarrow{HB}+\overrightarrow{AC}\cdot \overrightarrow{HC}\right)\\
			&=\dfrac{1}{4}\left[\overrightarrow{AB}\left(\overrightarrow{HC}+\overrightarrow{CB}\right)+\overrightarrow{AC}\left(\overrightarrow{HB}+\overrightarrow{BC}\right)\right]  \\
			&  =\dfrac{1}{4}\left(\overrightarrow{AB}\cdot \overrightarrow{HC}+\overrightarrow{AB}\cdot \overrightarrow{CB}+\overrightarrow{AC}\cdot \overrightarrow{HB}+\overrightarrow{AC}\cdot \overrightarrow{BC}\right)\\
			 &=\dfrac{1}{4}\left(\overrightarrow{AB}\cdot \overrightarrow{CB}+\overrightarrow{AC}\cdot \overrightarrow{BC}\right) \\
			& =\dfrac{1}{4}\left(\overrightarrow{AB}\cdot \overrightarrow{CB}-\overrightarrow{AC}\cdot  \overrightarrow{CB}\right)\\
			 &=\dfrac{1}{4}\overrightarrow{CB}\left(\overrightarrow{AB}-\overrightarrow{AC}\right)\\
			&=\dfrac{1}{4}\overrightarrow{CB}^2=\dfrac{1}{4}BC^2. 
		\end{align*}
		 Vậy $k=\dfrac{1}{4}=0{,}25$.
	 }
\end{bt}
